\chapter*{Abstract}
\addcontentsline{toc}{chapter}{Abstract}

Polymarket is the largest decentralised prediction market: contracts on real-world events
are traded on a central limit order book and settled on the Polygon blockchain, so that
every order-book update and every trade is publicly observable.  This thesis treats that
venue as a data-analysis and data-mining object and asks whether, and how, the direction
of its prices can be predicted over horizons from one minute to twelve hours.

The thesis first explains how a prediction market and Polymarket work, and introduces the
notions of time series, order books, graphs, learners and evaluation that it uses.  It then
follows twenty-one days of order-book, on-chain and market data from collection, through
preprocessing and audit, to a leakage-free three-class benchmark (down, unchanged, up), and
explores them as a price process, a transaction graph and a feature space.  The
methodology sets the objective and the criteria of comparison and presents five families
of approaches: causal rules on the price, stateless learners on hand-made features, a
continuous-time stateful model of the order-book stream, graph models on the
wallet--market transactions, and sources of information outside the book.  The
experiments compare these families on the same samples, with bootstrap intervals and a
locked test period, family by family and then across horizons, kinds of market and seeds.

The main findings are that a market's own activity is its most predictable feature; that
gradient-boosted trees on hand-made features are the strongest model from fifteen minutes
to twelve hours; that the stateful model of the book stream adds a small but reliable gain
at one minute, where the task is whether the price will move at all; and that the
transaction graph and the sources outside the book add no reliable gain over what the book
already shows.  The conclusions state the limits of the study, what a deployment would
need, and the future work.
